\section{Arithmetic and analytic preliminaries}
\label{sec:background}
\label{sec:preliminaries}

Write $K=\mathbb Q(\omega)$, $\mathcal O=\mathbb Z[\omega]$, and
$\lambda=\sqrt{-3}$, where $\omega=e^{2\pi i/3}$.
An element is \emph{primary} if it is congruent to $1$ modulo $3$.
Every nonzero ideal prime to $3$ has a unique primary generator.
Unless specified otherwise, sums over $a,b,c,d,r,u$ and their divisors
are over these generators. In contrast, Poisson frequencies range over
the full lattice indicated in the formula. We use
\[
 N(z)=|z|^2,\qquad \operatorname{Tr}(z)=z+\overline z,
 \qquad e(x)=\exp(2\pi i x).
\]
The functions $\mu$, $\tau$, and $\varphi$ are the ideal Mobius, divisor,
and Euler functions; thus $\varphi(r)=\#(\mathcal O/(r))^\times$.
The ideal von Mangoldt function is $\Lambda(n)=\log N\pi$ when
$n=\pi^k$ for a prime $\pi$ and $k\ge1$, and is zero otherwise.
We write $\Omega(n)$ for the number of prime factors counted with
multiplicity. The arithmetic function $\Lambda(n)$ is distinguished
by its argument from a completed $L$-function $\Lambda(s,\chi)$.


For a primary prime $\pi$, the cubic residue symbol
$\chi_\pi(v)$ is the element of $\{1,\omega,\omega^2\}$ congruent to
$v^{(N\pi-1)/3}$ modulo $\pi$ when $(v,\pi)=1$, and is zero otherwise.
Extend it multiplicatively in its lower argument. Put
\[
 G(b)=\frac{1}{\sqrt{Nb}}\sum_{v\bmod b}
       \chi_b(v)e\bigl(\operatorname{Tr}(v/b)\bigr),
 \qquad c_* =\frac{(2\pi)^{2/3}}{3\Gamma(2/3)}.
\]
Cubic reciprocity and its supplementary laws, in these conventions,
give $\chi_b(a)=\chi_a(b)$ for primary $a,b$.
The characters $b\mapsto\chi_b(\zeta\lambda^j)$, with $\zeta$ a unit,
have bounded conductor supported above $3$ and trivial infinite
parameter. We use these standard laws in the form recorded in
\cite[Section 2]{DR2024}.

Complete Gauss-sum evaluation and the Chinese remainder theorem give
\begin{equation}
 |G(a)|=\mu^2(a),\qquad
 G(ab)=G(a)G(b)\overline{\chi_b(a)}.
 \label{eq:mult}
\end{equation}
For completeness, the additive character is primitive away from $3$.
At a squarefree modulus the multiplicative character is primitive as
well, so the unnormalized sum has absolute value $\sqrt{Na}$.
At a repeated prime, summation over the final additive lift vanishes,
whereas the multiplicative character is inflated from a smaller modulus.
For coprime $a,b$, the two Chinese-remainder factors give
$\chi_b(a)\chi_a(b)=\overline{\chi_b(a)}$ by reciprocity.
If $(a,b)\ne1$, both sides of the second identity vanish.

For $h\ne0$ that is not a cube in $K$, the character
$b\mapsto\chi_b(h)$ is nonprincipal. Indeed it is the cubic character
of the nontrivial Kummer extension $K(\sqrt[3]{h})/K$.
The splitting criterion at unramified primes, together with the prime
ideal theorem for that extension, proves nontriviality.
Passing to the associated primitive character or adding coprimality
restrictions does not alter this assertion. Since $\mathcal O$ is
integrally closed, an element of $\mathcal O$ that is a cube in $K$
is already a cube in $\mathcal O$.

\subsection{Uniform conventions and classical inputs}

Throughout the proof, $X$ tends to infinity and $L=\log X$.
Lengths refer to norms. All dyads have fixed bounded endpoint ratios;
bounded lengths are allowed. All variables and reciprocal scale
parameters under consideration are bounded by fixed powers of $X$.
A sequence is \emph{divisor-bounded} if its absolute value is at most
$L^C\tau(n)^C$ for some fixed $C$.
Every fixed moment of this majorant on a norm dyad of length $B$ is
$O(BL^{C'})$, with $C'$ fixed. The same assertion holds after collecting
elements having the same rational norm: the number of ideals of norm
$m$, and their divisor multiplicities, are bounded by fixed powers of
the ordinary divisor function. These bounds follow from Euler products
and the usual mean estimates for fixed divisor functions.

The constants in $\ll_\epsilon$ may depend on all fixed support,
smoothness, and coefficient parameters. Each use of $X^\epsilon$ permits
an arbitrarily small positive exponent, chosen after those parameters.
We will explicitly retain logarithmic bounds when a small power loss
would not suffice.

We use Mellin inversion with the convention
\[
 \mathcal MW(s)=\int_0^\infty W(v)v^s\,\frac{dv}{v},\qquad
 W(v)=\frac{1}{2\pi i}\int_{(\sigma)}\mathcal MW(s)v^{-s}\,ds.
\]
Smooth functions in products or ratios of norm variables can consequently
be separated on fixed dyads. At fixed smoothness scale the Mellin
integrands decrease faster than every power of their imaginary
variables. A fixed number of derivatives bounded by powers of $L$
gives corresponding logarithmic costs. Derivative scales smaller than
one will be tracked when they are introduced.

The classical analytic inputs are the functional equation for primitive
Hecke characters, the prime ideal theorem and its Siegel--Walfisz form
over the fixed field $K$, and the ordinary mean-value inequality for
Dirichlet polynomials; see, for general analytic background,
\cite{IwaniecKowalski2004}. For the precise finite-order completion,
see \cite[Section 3.6]{Watkins2011}, and for entireness see
\cite[Section 2C]{ThornerZaman2019}. More precisely, if $\chi$ is primitive with
trivial infinite parameter and conductor norm $D$, then
\[
 \Lambda(s,\chi)
 =\left(\frac{\sqrt{3D}}{2\pi}\right)^s\Gamma(s)L(s,\chi),
 \qquad
 \Lambda(s,\chi)=\varepsilon_\chi\Lambda(1-s,\overline\chi),
 \quad |\varepsilon_\chi|=1.
\]
The field has discriminant $-3$ and one complex place. The complex
gamma factor in \cite{Watkins2011} is
$2(2\pi)^{-s}\Gamma(s)$; dividing the completion by the constant $2$
gives exactly this convention. For a nonprincipal character the
completed function is entire.
For every fixed $A,M>0$, nonprincipal cubic characters with primitive
conductor norm at most $(\log x)^A$ satisfy, uniformly,
\begin{equation}
 \sum_{N\pi\le x}\chi(\pi)\ll_{A,M}x(\log x)^{-M}.
 \label{eq:hecke-sw}
\end{equation}
Here is an explicit route to this uniformity. Apply
\cite[Theorem 1.4]{ThornerZaman2019} to the cyclic cubic extension of
$K$ cut out by $\chi$. Its three characters are
$1,\chi,\overline\chi$, and both nontrivial conductors have norm $D$.
The parameters of that theorem are $n_K=2$, $D_K=3$, and $Q=D$.
Weighting the three Frobenius-class counts by $\chi$ cancels their
principal terms. The possible exceptional character is real; in a
cyclic group of order three it must be trivial, so its term also
cancels. The resulting error is at most a constant times
\[
 \operatorname{Li}(x)\left\{
   \exp\!\left(-\frac{c\log x}{\log(12D)}\right)
       +\exp(-c\sqrt{\log x})\right\}+O(\log(2D)).
\]
For $D\le(\log x)^A$, the theorem's threshold
$x\ge(12D)^{c_2}$ holds eventually, and this error gives
\eqref{eq:hecke-sw} for every fixed $M$. The same theorem for a fixed
extension gives the prime ideal theorem and the nonprincipality test
used above. Standard partial summation gives the smooth and interval
versions of \eqref{eq:hecke-sw}. Any additional omitted Euler factors
will be accounted for through their deleted primes at the point of use.
For coefficients indexed by rational integers we use the classical
mean-value estimate of Montgomery and Vaughan~\cite{MontgomeryVaughan1974},
\begin{equation}
 \int_I\left|\sum_{n\le Z}v_n n^{it}\right|^2dt
 \ll (|I|+Z)\sum_{n\le Z}|v_n|^2,
 \label{eq:ordinary-mean-value}
\end{equation}
where $I$ is any interval. An arbitrary fixed translation of the height
is absorbed into the coefficients.

\subsection{Two large sieves}

\begin{theorem}[Heath-Brown's cubic large sieve]
\label{thm:cls}
Let $A,B\ge1$, and let $(u_b)$ be supported on primary squarefree
elements with $Nb\asymp B$. Then, for every $\epsilon>0$,
\[
 \sum_{Na\asymp A}\mu^2(a)
 \left|\sum_b u_b\chi_a(b)\right|^2
 \ll_\epsilon (AB)^\epsilon
 \bigl(A+B+(AB)^{2/3}\bigr)\sum_b|u_b|^2.
\]
The same assertion holds with upper cutoffs in place of dyads.
\end{theorem}

This is \cite[Theorem 2]{HB2000}. Its squarefree hypotheses apply to
both arguments, and will not be discarded without the following
explicit extensions.

The underlying planar large sieve is the number-field inequality
associated with Huxley~\cite{Huxley1968}; its exact $Z+Q^2$ scale
is also recorded in \cite[Equation (1.2)]{GaoZhao2022}. We give the
primitive-character argument needed here, including the absence of a
small-power conductor loss.

\begin{lemma}[The ordinary character large sieve]
\label{lem:ordinary-sieve}
For distinct primitive ray characters of trivial infinite parameter,
with conductor norm at most $Q$ and with conductor above $3$ dividing
a fixed modulus, and coefficients supported on a norm dyad of length
$Z$, one has
\[
 \sum_\chi\left|\sum_n v_n\chi(n)\right|^2
 \ll (Z+Q^2)\sum_n|v_n|^2.
\]
The implied constant is uniform when a fixed additional modulus above
$3$ is allowed. In particular, the bound applies to the cubic characters
indexed by coprime primary squarefree $p,q$ through
$b\mapsto\chi_b(pq^2)$, whose conductor away from $3$ is $pq$.
There is no $Q^\epsilon$ loss.
\end{lemma}

\begin{proof}
Fixing one of the finitely many classes at the bounded modulus above
$3$ reduces the claim to primitive group characters modulo a primary
$r$, with $Nr\ll Q$.
Let $S(y)=\sum_n v_n e(\operatorname{Tr}(ny))$.
Primitive Gauss-sum expansion and character orthogonality give
\[
 \sum_{\chi\bmod r}^{\mathrm{primitive}}
       \left|\sum_n v_n\chi(n)\right|^2
 \le \frac{\varphi(r)}{Nr}
       \sum_{\substack{h\bmod r\\(h,r)=1}}|S(h/r)|^2
 \le \sum_{\substack{h\bmod r\\(h,r)=1}}|S(h/r)|^2.
\]
Fixed factors arising from the trace-dual lattice are immaterial.
Distinct reduced fractions with denominator norms $O(Q)$ are separated
by $\gg Q^{-1}$ on the dual torus: their difference, after a lattice
translation, has a nonzero integral numerator and denominator of
absolute value $O(Q)$.

The planar additive large sieve for such a separated set has bound
$O(Z+Q^2)$. Here is the usual smoothing argument.
Take a nonnegative Schwartz majorant of the disk containing the support
of $(v_n)$, at scale $\sqrt Z$.
By duality and Poisson summation the resulting Gram matrix has entries
bounded by
\[
 C_N Z\sum_{\ell}(1+\sqrt Z\,|y-y'+\ell|)^{-N}.
\]
Packing a $Q^{-1}$-separated set in successive planar annuli bounds
each row sum by $O(Z+Q^2)$, for a fixed sufficiently large $N$.
The operator norm has the same bound, proving the assertion.
The exponents $1$ and $2$ at the primes dividing $pq$ distinguish the
cubic characters in the last assertion; bounded ramified factors only
introduce a bounded multiplicity.
\end{proof}

\begin{lemma}[Extensions of the cubic large sieve]
\label{lem:cubic-extensions}
The following estimates allow an arbitrarily small $X^\epsilon$ loss.
First, for divisor-bounded coefficients $(v_n)$ of norm length $Z$,
without a squarefree restriction,
\begin{equation}
 \sum_{Na\asymp P}\mu^2(a)
       \left|\sum_n v_n\chi_a(n)\right|^2
 \ll_\epsilon X^\epsilon Z
       \bigl(P+Z+(PZ)^{2/3}\bigr).
 \label{eq:inner-extension}
\end{equation}
Second, for coefficients $(u_b)$ on primary squarefree elements of
norm length $P$,
\begin{equation}
 \sum_{Na\asymp J}\left|\sum_b u_b\chi_b(a)\right|^2
 \ll_\epsilon X^\epsilon
       \bigl(J+(JP)^{2/3}+J^{1/3}P\bigr)\sum_b|u_b|^2.
 \label{eq:outer-extension}
\end{equation}
The latter also holds when $a$ ranges over all nonzero elements of
$\mathcal O$ of norm $\asymp J$. If $J\gg P$, its last term is
absorbed by $(JP)^{2/3}$.
\end{lemma}

\begin{proof}
For the first estimate write uniquely $n=s t^2 c^3$, with $s,t$
coprime and squarefree. Restrict their norms to dyads $S,V,C$, so that
$SV^2C^3\asymp Z$.
For each fixed $t,c$, the factors involving them have modulus at most
one on the outer variable. Apply Theorem~\ref{thm:cls} to the $s$ sum
and then use the triangle inequality in the outer squared-sum norm.
The squared norm of the restricted coefficient sequence is
$O_\epsilon(X^\epsilon S)$.
There are $O(VC)$ choices of $t,c$, and the resulting bound is
\begin{equation}
 X^\epsilon (VC)^2S
       \bigl(P+S+(PS)^{2/3}\bigr).
 \label{eq:inner-extension-dyad}
\end{equation}
Since $(VC)^2S\ll Z$ and $S\ll Z$, summing the logarithmically many
dyads proves \eqref{eq:inner-extension}. We retain
\eqref{eq:inner-extension-dyad} when a restriction makes $S$ shorter.

For the second estimate factor the outer argument in the same way.
Fixing $t,c$ when $S\ge V$, and applying Theorem~\ref{thm:cls} to $s$,
costs
\[
 X^\epsilon VC\bigl(S+P+(SP)^{2/3}\bigr)\sum_b|u_b|^2.
\]
Fixing $s,c$ when $S<V$ gives instead
\[
 X^\epsilon SC\bigl(V+P+(VP)^{2/3}\bigr)\sum_b|u_b|^2,
\]
with character conjugation if necessary.
Using $SV^2C^3\ll J$ and the indicated order of $S,V$, each expression
is bounded by the right side of \eqref{eq:outer-extension}.
The coprimality restriction imposed by $c$ is simply a restriction of
the coefficients for each fixed $c$.
For a full-lattice argument write
$a=\zeta\lambda^i s t^2 c^3$, $i\in\{0,1,2\}$, allowing powers of
$\lambda$ in $c$. The finitely many unit and residual ramified factors
are coefficient twists. The same proof applies.
\end{proof}

\subsection{The unconditional level input}

We use only the unconditional summation formula of Dunn and Radziwill,
not their subsequent estimates that assume GRH. To identify the exact
input, equation numbers in the following statement refer to
\cite[arXiv version 3, Proposition 5.3]{DR2024}.

\begin{theorem}[Level Voronoi formula, the required form]
\label{thm:voronoi-background}
Let $r$ be primary and squarefree, and let $W$ be smooth and compactly
supported in $(0,\infty)$. Put $R=Nr$ and
$A_0=(2\pi)^{5/3}/(3^{7/2}\Gamma(2/3))$.
The completed sum
\[
 \mathcal C_r(V;W)=
 \sum_{\substack{d,u\\(d,r)=1}}|d|G(ur)
             W\left(\frac{N(ud^3)}{V}\right)
\]
has main term
$A_0\mathcal MW(5/6)V^{5/6}\varphi(r)R^{-7/6}$.
Its remaining term is
\[
 \frac{R^{1/2}}{3^{7/2}(2\pi)^2}
 \sum_{\substack{0\ne\nu\in\lambda^{-1}\mathcal O, d\\(d,r)=1}}
 \frac{a_r(\nu)b_r(\nu)}{N\nu\,(Nd)^{5/2}}
 \check W\left(\frac{(2\pi)^4N(d^3\nu)V}{R^2}\right),
\]
where, for sufficiently small $\sigma>0$,
\[
 \check W(v)=\frac{1}{2\pi i}\int_{(-\sigma)}v^z
 \frac{\Gamma(5/6-z)\Gamma(7/6-z)}
      {\Gamma(z-1/6)\Gamma(z+1/6)}\mathcal MW(z)\,dz.
\]
The coefficients depend only on $r$ and $\nu$, independently of $V$
and $W$. They vanish outside representations
\[
 \nu=\lambda^j\zeta h w(h')^3,\quad j\ge-1,\quad
 h,h'\mid r^\infty,\quad (w,r)=1,\quad \mu^2(hw)=1,
\]
where $h,h',w$ are primary and $\zeta$ is a unit.
On this support, they are
\[
 a_r(\nu)=a^\star(\lambda^{-3}\nu),\qquad
 b_r(\nu)=\mu\left(\frac{r}{(\lambda\nu,r)}\right)
          \frac{\varphi(r)}{\varphi\left(r/(\lambda\nu,r)\right)},
\]
where $a^\star$ is the theta coefficient in
\cite[Equation (5.74)]{DR2024}; the second formula is its
Equation (5.68). The index $\lambda^{-3}\nu$ is the one supplied
by its coefficient identity (5.79).
They satisfy
\[
 |a_r(\nu)|\ll 3^{\max(j,0)/3}|h'|,
 \qquad |b_r(\nu)|\le N((\lambda\nu,r)).
\]
In particular, if $V,V^{-1},R\ll X^{C_0}$ for fixed $C_0$, $W$ is
supported in a fixed compact subinterval of $(0,\infty)$, and each
fixed derivative bound is at most a fixed logarithmic power, then
\begin{equation}
 \mathcal C_r(V;W)=
 A_0\mathcal MW(5/6)V^{5/6}\frac{\varphi(r)}{R^{7/6}}
 +O_\epsilon(X^\epsilon R^{1/2}).
 \label{eq:voronoi}
\end{equation}
\end{theorem}

\begin{proof}[Derivation of the stated consequence]
The summation formula and coefficient information are the cited
background theorem, with Equations (5.7), (5.13)--(5.14),
(5.74), and (5.79)--(5.81). In particular, the dual theta coefficient
is read from the proof formula (5.79), where its index is
$\lambda^{-3}\nu$; this avoids the repeated ramified exponent in
the printed (5.67). The theta formulas give
$|a_r(\nu)|\ll3^{\max(j,0)/6}|h'|$, which implies the stated
weaker bound; the finitely many exceptional ramified cases have bounded
factors. Multiplying the original summation formula by $G(r)$ uses
$G(r)\overline{g(r)}=R^{1/2}$, where $g(r)=\sqrt R\,G(r)$.
This proves the displayed normalization of both terms.

We give the absolute convergence argument for \eqref{eq:voronoi}.
On $\operatorname{Re}z=-\sigma$, rapid Mellin decay and Stirling's
formula bound the transform by
\[
 X^{O(\sigma)}L^{O(1)}
              \bigl(N(d^3\nu)\bigr)^{-\sigma}.
\]
For a prime of norm $p$ dividing $r$, the exponents in $h,h'$ are
$e\in\{0,1\}$ and $k\ge0$. Its absolute local sum is bounded by
\[
 \sum_{e=0}^1\sum_{k\ge0}
 p^{1_{e+k>0}}p^{k/2}p^{-(e+3k)(1+\sigma)}
 =1+p^{-\sigma}+O_\sigma(p^{-3/2-3\sigma}).
\]
For every fixed $\eta>0$, the product over $p\mid r$ is
$O_{\sigma,\eta}(R^\eta)$: separate the finitely many small primes,
then bound each remaining factor by $p^\eta$.
The free $w$ sum converges at $1+\sigma$, the ramified sum is dominated
by $\sum_{j\ge0}3^{-j(2/3+\sigma)}$, and the $d$ sum converges at
$5/2+3\sigma$.
Choosing $\sigma$ and $\eta$ sufficiently small in terms of the final
$\epsilon$ proves \eqref{eq:voronoi}.
\end{proof}

\begin{theorem}[The metaplectic mean-square input]
\label{thm:hb-height}
For primary squarefree $r$, the Dirichlet series
\[
 f_r(s)=\sum_u G(ru)(Nu)^{-s},\qquad \operatorname{Re}s>1,
\]
continues meromorphically to $\operatorname{Re}s>1/2$, with at most
a simple pole at $5/6$. It has polynomial growth in closed strips
there, away from its pole. For sufficiently small $\epsilon>0$ and
every $T\ge1$,
\begin{equation}
 \int_{-T}^T|f_r(1/2+\epsilon+it)|^2dt
 \ll_\epsilon (Nr)^{1/2+4\epsilon}T^2.
 \label{eq:hb-mean-square}
\end{equation}
Its residue $\rho_r$ satisfies $|\rho_r|\ll (Nr)^{-1/6}$.
\end{theorem}

The mean-square assertion is \cite[Section 3, Lemma 3]{HB2000}, with
angular index zero; the series there is exactly $f_r$ in this
normalization. We spell out the other analytic properties separately.
Multiplicativity identifies $f_r$ with $G(r)$ times the series in
\cite[Proposition 5.2]{DR2024}. Its completion includes the factor
$\sum_{(d,r)=1}(Nd)^{-(3s-1/2)}$, whose Euler product is nonzero for
$\Rea s>1/2$. Thus division introduces no further pole in this
half-plane. The residue is

\begin{equation}
 \rho_r=A_0\frac{\varphi(r)}{(Nr)^{7/6}}
       \left(\sum_{(d,r)=1}(Nd)^{-2}\right)^{-1}.
 \label{eq:typeI-residue}
\end{equation}
The last sum is at least one, proving the bound without a small power
loss in the level.

For the growth needed in a contour shift, put $R=Nr$. In the notation
of \cite[Section 3, Equation (14)]{HB2000}, write
$Z(q,z)=\zeta_K(3z-2;1_3)\psi(q,z)$, where $1_3$ denotes omission
of the prime above $3$. Equation (19) there bounds this function at
large height by
\[
 |Z(q,\alpha+iy)|\ll_e
 (Nq)^{(3/2+e-\alpha)/2}(1+y^2)^{3/2+e-\alpha},
 \qquad \tfrac12\le\alpha\le\tfrac32.
\]
Only $|y|\ge1$ is used, away from the possible pole. The finite-divisor
identity at the end of that section gives
$f_r(s)=R^s\psi_r(1,s+1/2)$ and, for $\Rea z\ge1+\delta$,
\[
 |\psi_r(1,z)|\ll_\delta R^{-1/2}
       \sum_{d\mid r}(Nd)^{-1/2}|\psi(r/d,z)|.
\]
The inverse zeta factor and the finite Euler factors in this identity
are bounded by absolutely convergent products in that region.
Consequently, for $1/2+\delta\le\sigma\le.9$, $|t|\ge1$, and any
fixed $e,\nu>0$,
\[
 |f_r(\sigma+it)|\ll_{\delta,e,\nu}
 R^{\sigma/2+e/2+\nu}(1+t^2)^{1+e-\sigma}.
\]
Here the divisor sum costs at most $\tau(r)\ll_\nu R^\nu$.

To cover the rest of a strip, set
$S_r(U)=\sum_{Nu\le U}G(ru)$. For $U\ge R^2$,
\cite[Lemma 1]{HB2000}, with its parameter $1/12$, applies because
$R\le(RU)^{1/3}$ and gives
\[
 |S_r(U)|\ll R^{-1/6}U^{5/6}+R^{1/4}U^{3/4}
           \ll R^{1/3}U^{5/6}.
\]
For $1\le U<R^2$ the same final bound follows from the trivial bound
$O(U)$. Abel summation now gives, for $\Rea s>5/6$,
\[
 f_r(s)=s\int_1^\infty S_r(y)y^{-s-1}\,dy,
 \qquad
 |f_r(s)|\ll R^{1/3}\frac{|s|}{\Rea s-5/6}.
\]
This proves polynomial growth on $\Rea s\ge.9$, overlapping the
previous range. Bounded heights away from the pole follow from
meromorphy for each fixed level. These bounds justify the contour
shifts below; their quantitative height estimate is
\eqref{eq:hb-mean-square}.


\section{Type I estimates}
\label{sec:typeI}

The Type I sums have arbitrary coefficients in the level $r$
and a smooth unrestricted sum over $u$. The level formula evaluates
this inner sum after completion by cubes. We remove that completion by
M\"obius inversion: its residue gives the squarefree model, while the
discarded cube divisors are bounded by counting at small levels and by
the cubic large sieve at the remaining levels. At large Mellin heights
we instead estimate the Gauss-sum term directly from the metaplectic
mean square.

\begin{proposition}[Type I comparison]
\label{prop:typeI}
Suppose $RU\asymp X$ and $1\le R\ll X^{.51}$.
Let $(a_r)$ be divisor-bounded and supported on $Nr\asymp R$.
Let $W_r$ have support in a fixed compact subinterval of $(0,\infty)$,
with each fixed derivative bounded by a fixed power of $L$, uniformly
in $r$. Then
\[
 \sum_r a_r\sum_u
 \left(G(ru)-c_*\frac{\mu^2(ru)}{N(ru)^{1/6}}\right)W_r(Nu/U)
 \ll X^{5/6-\eta}
\]
for some fixed $\eta>0$. One may take $\eta=1/100$, with constants
depending on the stated fixed parameters.
\end{proposition}

\begin{proof}
Nonsquarefree $r$ contribute zero to both terms. For squarefree $r$,
invert cube completion by the exact identity
\begin{equation}
 \sum_uG(ru)W_r(Nu/U)
 =\sum_{(c,r)=1}\mu(c)|c|\,
       \mathcal C_r(U/(Nc)^3;W_r).
 \label{eq:inverse-completion}
\end{equation}
Indeed the coefficient at a cube $e^3$ is
$|e|\sum_{c\mid e}\mu(c)$, equal to zero unless $e=1$.
All sums in this identity are finite at the given compact support.
Use \eqref{eq:voronoi} for $Nc\le C_*$, where
\[
 C_*=
 \begin{cases}
 X^{.13},& R\le X^{.40},\\
 X^{.02},& R>X^{.40}.
 \end{cases}
\]
After summing $r$, the total error is
\begin{equation}
 \ll X^\epsilon R^{3/2}C_*^{3/2}\ll X^{.795+\epsilon}.
 \label{eq:short-completion-error}
\end{equation}
The main term, completed to all $c$, is
\begin{equation}
 A_0\mathcal MW_r(5/6)U^{5/6}
 \frac{\varphi(r)}{(Nr)^{7/6}}
 \sum_{(c,r)=1}\frac{\mu(c)}{(Nc)^2}.
 \label{eq:typeI-main}
\end{equation}
Its omitted tail, after summing $r$, is
$O(X^{5/6+\epsilon}C_*^{-1})$.

We compare \eqref{eq:typeI-main} with the model directly.
The coset $1+3\mathcal O$ has Euclidean area density $2/(9\sqrt3)$,
hence radial density $2\pi/(9\sqrt3)$ in the norm variable.
Square-divisor and coprimality inversion show, for a log-smooth
weight $W$, that
\[
 \sum_{(u,r)=1}\mu^2(u)W(Nu/U)
 =\frac{2\pi U}{9\sqrt3}\mathcal MW(1)
       \frac{\varphi(r)}{Nr}
       \sum_{(c,r)=1}\frac{\mu(c)}{(Nc)^2}
       +O_\epsilon(X^\epsilon\sqrt U).
\]
To justify the error, truncate square divisors at $Nc\ll\sqrt U$.
Each relevant sublattice count differs from its area by at most a
constant times one plus its scaled radius. Summing these errors and
the divisors of $r$ costs only $X^\epsilon\sqrt U$; completing the
area term has the same error.
Apply this formula to $v^{-1/6}W_r(v)$ and use
\[
 c_*\frac{2\pi}{9\sqrt3}=A_0.
\]
The model therefore equals \eqref{eq:typeI-main}, with total error
over $r$ bounded by
\begin{equation}
 X^\epsilon R^{5/6}U^{1/3}
 \asymp X^{1/3+\epsilon}R^{1/2}
 \ll X^{.589+\epsilon}.
 \label{eq:typeI-lattice-error}
\end{equation}

It remains to bound the omitted part of \eqref{eq:inverse-completion}.
Combine $c,d$ into $e=cd$, and restrict $Ne$ to a dyad of length
$E\gg C_*$. Its coefficient is divisor-bounded times $E^{1/2}$;
the restriction $(e,r)=1$ may be kept in the outer coefficient.
Trivial counting gives
\begin{equation}
 \ll X^{1+\epsilon}E^{-3/2}.
 \label{eq:typeI-trivial-tail}
\end{equation}
When $R\le X^{.40}$ this is $O(X^{.805+\epsilon})$.

For the other range fix $e$ and put $U_e=U/E^3$.
Using \eqref{eq:mult}, Cauchy in $r$, and
Theorem~\ref{thm:cls}, the bilinear sum in $r,u$ costs
\[
 \ll X^\epsilon\sqrt{RU_e}
       \bigl(R+U_e+(RU_e)^{2/3}\bigr)^{1/2}.
\]
The automatic vanishing restricts both variables to squarefree elements.
Mellin inversion separates the weights; the arbitrary dependence of
$W_r$ on $r$ changes only the outer coefficients and a common
integrable majorant. Summing the absolute $e$ coefficients on the dyad
costs $O(X^\epsilon E^{3/2})$. Thus the tail on that dyad is
\begin{equation}
 \ll X^{1/2+\epsilon}
       \left(R+U/E^3+(X/E^3)^{2/3}\right)^{1/2}.
 \label{eq:typeI-sieve-tail}
\end{equation}
For $X^{.40}<R\ll X^{.51}$ and $E\gg X^{.02}$, its three terms
have exponents at most
\[
 .755,\qquad .77,\qquad
 \frac12+\frac{1-.06}{3}=\frac56-\frac1{50}.
\]
The same last upper exponent bounds the omitted main-term tail in
this range. Equations \eqref{eq:short-completion-error}--
\eqref{eq:typeI-sieve-tail}, with logarithmically many dyads and
sufficiently small epsilon losses, prove the assertion with
$\eta=1/100$.
\end{proof}

\begin{proposition}[Type I at large Mellin heights]
\label{prop:typeIheight}
Suppose $RU\asymp X$, and let $(a_r)$ be divisor-bounded on
$Nr\asymp R$. Suppose the smooth compactly supported weights $W_r$
have fixed support ratios and uniformly bounded derivatives,
independently of $X$. For $T\ge2$ and every fixed $D>0$,
\begin{align}
 &\int_{|t|\asymp T}\sum_{Nr\asymp R}|a_r|
       \left|\sum_uG(ru)(Nu)^{it}W_r(Nu/U)\right|\frac{dt}{T}
 \notag\\
 &\hspace{12mm}\ll_\epsilon
       X^{1/2+\epsilon}R^{3/4}T^{1/2}
       +O_D\bigl(X^{5/6}L^{C}T^{-D}\bigr),
 \label{eq:typeI-height}
\end{align}
where $C$ depends only on the coefficient bounds.
\end{proposition}

\begin{proof}
Again only squarefree $r$ contribute.
For each such $r$, Mellin inversion on a line to the right of one gives
\[
 \sum_uG(ru)(Nu)^{it}W_r(Nu/U)
 =\frac{1}{2\pi i}\int_{(2)}
       \mathcal MW_r(s)U^s f_r(s-it)\,ds.
\]
Shift to $\sigma=1/2+\epsilon'$, with $\epsilon'>0$ small.
Theorem~\ref{thm:hb-height} and rapid Mellin decay justify this shift.
The residue is
\[
 \rho_r U^{5/6+it}\mathcal MW_r(5/6+it)
 \ll_D U^{5/6}R^{-1/6}(1+|t|)^{-D}.
\]
For the new integral use Minkowski and Cauchy in $t$.
For every real $v$, \eqref{eq:hb-mean-square}, applied on a symmetric
interval of radius $O(T+|v|+1)$, gives
\[
 \left(\int_{|t|\asymp T}
       |f_r(\sigma+i(v-t))|^2\frac{dt}{T}\right)^{1/2}
 \ll_{\epsilon'}R^{1/4+2\epsilon'}
              \frac{T+|v|+1}{\sqrt T}.
\]
Since $\mathcal MW_r(\sigma+iv)$ decreases rapidly, uniformly in $r$,
the mean absolute value of the new integral is
\[
 \ll_{\epsilon'} U^{1/2+\epsilon'}R^{1/4+2\epsilon'}\sqrt T.
\]
Finally $\sum_{Nr\asymp R}|a_r|\ll RL^C$.
Summing the last bound and the residues, and taking $\epsilon'$
sufficiently small in terms of $\epsilon$, proves
\eqref{eq:typeI-height}.
\end{proof}

The two applications of Proposition~\ref{prop:typeIheight} are
\[
 R\le X^{.40},\quad T\ll X^{.01},
 \qquad\text{and}\qquad
 R\le X^{1/3-\kappa/2},\quad T\ll X^{1/6+\rho}.
\]
The respective first-term exponents are at most $.805+\epsilon$ and
$5/6-3\kappa/8+\rho/2+\epsilon$.
Both have a fixed power gap when $\rho$ is sufficiently small in terms
of $\kappa$. Above a sufficiently large fixed power of $L$, the pole
term has any prescribed logarithmic saving by choosing $D$ large.
